%==============================================================================
\chapter{Conclusion and Future Scope}
\label{ch:conclusion}
%==============================================================================

\section{Conclusion}
\label{sec:conclusion}

This thesis presented TutorAI, a complete, working system that turns a typed
topic into a short, self-explaining lesson: a vector diagram whose parts
carry stable semantic identifiers, segment-level narration tied to those
identifiers, and per-segment synthesized audio that drives synchronized
highlighting on an Android device --- fully offline once downloaded, with an
on-device history and no user accounts.

All six objectives set out in Section~\ref{sec:objectives} were met. An
agentic generation pipeline was designed and implemented as a LangGraph state
graph that transforms a free-text topic into a validated lesson
(Objective~1). A deterministic, non-LLM validation gate with bounded
critique-and-refine and repair loops guarantees that no lesson with a broken
narration reference is ever cached or served, and this guarantee was verified
in both directions --- rejection and repair of broken lessons, and clean
failure when repair is exhausted (Objective~2). A stateless FastAPI backend
provides the asynchronous job model, the shared topic-keyed cache with
request coalescing, the asset store, and swappable Gemini provider adapters
(Objective~3). An offline-first Android client built with Jetpack Compose,
Clean Architecture, and MVVM generates lessons via job polling, plays them
with deterministic segment-synchronized highlighting, and persists them for
network-free replay (Objective~4). The engineering practice throughout ---
UML-first design, classical patterns at the seams, 48 hermetic key-free
tests, containerization, and CI/CD for both backend and client --- makes the
system reproducible from a fresh clone (Objective~5). Finally, the realized
system was verified end to end against live models, producing consistent
lessons across topics as different as photosynthesis, the Pythagorean
theorem, and the Transformer architecture (Objective~6).

The central insight of the work is architectural rather than about any
single model. By splitting generation into a small graph of steps, adding a
bounded critique-and-refine loop for the drawing and a bounded repair loop
for the narration, and placing a deterministic, non-LLM validator in front
of the cache, a probabilistic generator is made reliable enough to power a
feature that depends on exact part-to-word correspondence. Around this core,
a stateless caching backend keeps cost down, provider adapters keep the
vendor choice soft, and an offline-first client keeps the lesson usable
without a network. The pattern --- probabilistic quality mechanisms in front
of a deterministic correctness gate --- transfers directly to any
application whose LLM outputs must satisfy machine-checkable invariants.

\section{Future Scope}
\label{sec:futurescope}

Future work falls into three groups; none changes the central design ---
each extends it.

\textbf{Evaluation.} The most important next step is empirical: a controlled
learning study comparing TutorAI lessons against static diagrams and plain
text, grounded in the multimedia-learning predictions that motivated the
product form; a larger automatic study of diagram quality and validator
pass/repair/fail rates across hundreds of topics; and on-device latency,
memory, and energy measurements across a range of phones.

\textbf{Capability.} Several extensions are natural. Optional word-level
highlighting where TTS timing metadata is available, degrading gracefully to
segment-level sync offline. Richer diagram styles and layout hints from the
planner. A research or planning step in front of the drawing (for example, a
retrieval or deep-research agent behind the existing
\texttt{LLMProvider.plan\_concepts} port) for harder or fast-moving topics.
Multi-language lessons --- the manifest already carries language and voice
parameters. Accessibility refinements such as adjustable playback speed and
transcript export.

\textbf{Scale.} The development profile's in-process worker and filesystem
store can be promoted to a separate worker pool, a managed PostgreSQL cache,
and cloud object storage behind the same repository interfaces, with no
call-site changes. Server-sent events can replace job polling. A lightweight
API key can be introduced through the header slot the contract already
reserves, if usage ever warrants it.
